\section{Requisitos de componentes de inteligencia artificial}

Especificar un componente de aprendizaje autom\'atico no es especificar una funci\'on determinista: el comportamiento se induce de datos, la salida es probabil\'istica y la calidad depende tanto del modelo como del conjunto de datos y del contexto de uso. Por eso la especificaci\'on debe fijar medidas de desempe\~no, requisitos sobre los datos y requisitos de calidad que no aparecen en el software convencional. Entre ellos, la \textbf{explicabilidad} se ha consolidado como requisito no funcional de pleno derecho, con definici\'on, modelo y cat\'alogo propios \cite{chazette2020explainability,chazette2021exploring}, y con criterios para evaluarla a lo largo del ciclo de vida \cite{chazette2022explainable}. Ese es el marco que sustenta el RNF-10 ya presente en el ERS.

El sistema tiene dos componentes basados en aprendizaje autom\'atico, ambos ya identificados en el ERS como RF-09 y RF-14 (\S3.2 del ERS, \S\ref{sec:ers-final}), pero sin especificaci\'on de rendimiento, equidad ni explicabilidad m\'as all\'a de RNF-10 (explicabilidad de RF-09, ya existente). Esta secci\'on formaliza ambos componentes con los cinco bloques exigidos por la gu\'ia PE5 y cubre los m\'inimos por componente: 3 requisitos funcionales, 3 no funcionales de rendimiento, 2 de equidad y 1 de explicabilidad. En total se a\~naden seis RF de IA (RF-IA-01 a RF-IA-06) y once RNF de IA (RNF-IA-01 a RNF-IA-11). Ninguno queda aislado: todos refinan RF-09 o RF-14, ya especificados en el ERS, y cada uno tiene su caso de prueba conceptual en \texttt{04\_Trazabilidad/casos\_prueba\_conceptuales.csv}.

\subsection{IA-01: Predictor de fallas de equipos (RF-09)}
\textbf{Tarea y tipo:} clasificaci\'on binaria (riesgo de falla alto/bajo) por equipo y ventana temporal, sobre series de tiempo de sensor. \textbf{Sustenta:} RF-09, RNF-10 (explicabilidad, ya existente en el ERS).

\textbf{(a) Descripci\'on del componente.} Entradas: temperatura, humedad, ciclos de encendido/apagado del equipo, alertas previas (RF-08/RF-10). Salida: puntaje 0--1, etiqueta alto/bajo riesgo, variables m\'as influyentes. Consumidor: Personal de Infraestructura y Mantenimiento, v\'ia alerta predictiva (CU-16). \emph{Fallback:} una predicci\'on de alto riesgo no ejecuta ninguna acci\'on autom\'atica; genera una alerta informativa que un humano eval\'ua.

\textbf{Requisitos funcionales del componente.} Ninguno queda aislado: los tres refinan RF-09, ya especificado en el ERS.

\begingroup
\small
\begin{longtable}{|p{2.1cm}|p{9.4cm}|p{3.4cm}|}
\hline
\cellcolor{sigagreen}\textcolor{white}{\textbf{ID}} & \cellcolor{sigagreen}\textcolor{white}{\textbf{Enunciado verificable}} & \cellcolor{sigagreen}\textcolor{white}{\textbf{Traza}} \\ \hline
RF-IA-01 & El sistema debe calcular, para cada equipo con al menos 90 d\'ias de historial continuo, un puntaje de riesgo de falla en el rango 0--1 referido a una ventana de los pr\'oximos 7 d\'ias, y recalcularlo al menos una vez al d\'ia. & Refina RF-09; verifica CP-IA-01; medido por RNF-IA-01 y RNF-IA-02 \\ \hline
RF-IA-02 & Cuando el puntaje de un equipo supere el umbral de alto riesgo, el sistema debe generar una alerta predictiva de prioridad Media asociada a ese equipo y aula, visible en la vista de alertas, \textbf{sin} crear autom\'aticamente un ticket de mantenimiento. & Refina RF-09; se apoya en RF-08 y CU-16; verifica CP-IA-02 \\ \hline
RF-IA-03 & Para cada predicci\'on de alto riesgo mostrada, el sistema debe presentar las variables de sensor que m\'as contribuyeron a esa predicci\'on, en lenguaje natural. & Refina RF-09; operacionaliza RNF-10; verifica CP-IA-03 \\ \hline
\caption{Requisitos funcionales del componente IA-01}
\end{longtable}
\endgroup

\textbf{(b) Datos de entrenamiento.} Origen: historial acumulado por RF-01/RF-03/RF-05 y RF-10. Volumen m\'inimo: 90 d\'ias continuos por equipo (consistente con SUP-06). Etiquetado autom\'atico desde tickets correctivos cerrados (RF-12). Calidad m\'inima: $\geq$80\% de lecturas presentes. Sesgo conocido: equipos de bajo uso reciente subrepresentados: se declara, no se corrige artificialmente. No es dato personal; no aplica LOPDP directamente.

\textbf{(c) M\'etricas de \'exito.}
\begin{itemize}[nosep]
\item RNF-IA-01: F1 $\geq$0,70 en conjunto de prueba (holdout temporal) antes de cada despliegue.
\item RNF-IA-02: latencia de inferencia $\leq$2~s (percentil 95).
\item RNF-IA-03: reevaluaci\'on cada 30 d\'ias; retiro autom\'atico si F1 $<$0,60 en dos evaluaciones consecutivas.
\end{itemize}

\textbf{(d) Requisitos \'eticos y de privacidad.}
\begin{itemize}[nosep]
\item RNF-IA-04 (equidad): diferencia de falsos negativos $\leq$10 puntos porcentuales entre aulas de uso alto y bajo.
\item RNF-IA-05 (equidad): diferencia de falsos negativos $\leq$10 puntos porcentuales entre edificios/bloques.
\item RNF-10 (explicabilidad, reutilizado): explicaci\'on en $\leq$2~s, $\leq$60 palabras, comprensi\'on $\geq$4/5 Likert.
\item Finalidad: mantenimiento preventivo, no evaluaci\'on de docentes. Sin datos personales. Supervisi\'on humana obligatoria antes de crear un ticket.
\end{itemize}

\textbf{(e) Plan de monitoreo.} Indicadores: F1 mensual, falsos negativos por segmento, latencia p95. Umbral de deriva: ca\'ida de F1 $>$15 puntos respecto a la l\'inea base. Reentrenamiento: cada 90 d\'ias o al activarse el umbral.


El componente no se especifica en el vac\'io: opera dentro de CU-16, cuyo diagrama de secuencia (Figura~\ref{fig:seq-uc16}) fija en qu\'e orden colaboran los objetos y, con ello, d\'onde se inserta la inferencia del modelo dentro del flujo. Esa ubicaci\'on es la que hace medible el umbral de latencia de RNF-IA-02: el l\'imite de dos segundos se aplica al paso de inferencia (mensaje 1.1.8), no al caso de uso completo. Conviene precisar una diferencia de granularidad entre el diagrama y RF-IA-01, para que no se lea como contradicci\'on: el diagrama, elaborado en la Entrega 3, muestra el indicador que se presenta al usuario, categ\'orico (\emph{Low}, \emph{Medium}, \emph{High}); RF-IA-01 especifica la magnitud que el modelo calcula, continua en el rango 0--1. La categor\'ia es la discretizaci\'on de ese puntaje para la interfaz, y los umbrales de corte quedan como decisi\'on de dise\~no pendiente, no como requisito ya fijado.

\begin{figure}[H]
\centering
\makebox[\textwidth][c]{\includegraphics[width=1.25\textwidth]{seq_UC16_predict_equipment_failures.png}}
\caption{Diagrama de secuencia de CU-16, predicci\'on de fallas de equipos}
\label{fig:seq-uc16}
\end{figure}

\subsection{IA-02: An\'alisis de patrones de consumo energ\'etico (RF-14)}
\textbf{Tarea y tipo:} detecci\'on de anomal\'ias + agrupamiento no supervisado sobre series de consumo. \textbf{Sustenta:} RF-14, RF-17 (consumidor del resultado).

\textbf{(a) Descripci\'on del componente.} Entradas: ocupaci\'on (RF-03), estado de climatizaci\'on/proyector (RF-04/RF-05), lectura directa de consumo cuando exista sensor dedicado. Salida: lista de anomal\'ias de consumo (aula, franja horaria, magnitud), integrada al reporte administrativo (RF-17/RF-18). Consumidor: Autoridades de la Facultad, Personal Administrativo. \emph{Fallback:} una anomal\'ia mal detectada solo agrega una l\'inea a un reporte revisado por un humano; no ejecuta ninguna acci\'on sobre equipos.

\textbf{Requisitos funcionales del componente.} Los tres refinan RF-14 y alimentan RF-17; ninguno queda aislado.

\begingroup
\small
\begin{longtable}{|p{2.1cm}|p{9.4cm}|p{3.4cm}|}
\hline
\cellcolor{sigagreen}\textcolor{white}{\textbf{ID}} & \cellcolor{sigagreen}\textcolor{white}{\textbf{Enunciado verificable}} & \cellcolor{sigagreen}\textcolor{white}{\textbf{Traza}} \\ \hline
RF-IA-04 & El sistema debe calcular, para cada aula con al menos 60 d\'ias de historial, una l\'inea base de consumo por franja horaria y d\'ia de la semana, e indicar expl\'icitamente si esa l\'inea base proviene de un sensor de consumo real o del indicador indirecto (ocupaci\'on y estado de equipos). & Refina RF-14; se apoya en RF-03/RF-04/RF-05; verifica CP-IA-04 \\ \hline
RF-IA-05 & El sistema debe ejecutar, una vez al d\'ia, la detecci\'on de desviaciones respecto a la l\'inea base y producir la lista de anomal\'ias con aula, franja horaria y magnitud de la desviaci\'on. & Refina RF-14; verifica CP-IA-05; medido por RNF-IA-06, RNF-IA-07 y RNF-IA-08 \\ \hline
RF-IA-06 & El sistema debe incorporar las anomal\'ias detectadas al reporte administrativo del periodo, agregadas por aula y franja horaria, indicando para cada una la variable que m\'as contribuy\'o a la desviaci\'on. & Refina RF-14; alimenta RF-17/RF-18; operacionaliza RNF-IA-10 y RNF-IA-11; verifica CP-IA-06 \\ \hline
\caption{Requisitos funcionales del componente IA-02}
\end{longtable}
\endgroup

\textbf{(b) Datos de entrenamiento.} Volumen m\'inimo: 60 d\'ias por aula. No requiere etiquetado manual (modelo no supervisado). Sesgo conocido: las aulas sin sensor de consumo directo dependen de un indicador indirecto menos preciso: se declara expl\'icitamente en el reporte. Dato agregado de aula, no personal.

\textbf{(c) M\'etricas de \'exito.}
\begin{itemize}[nosep]
\item RNF-IA-06: tasa de falsos positivos $\leq$15\% sobre validaci\'on manual trimestral.
\item RNF-IA-07: an\'alisis por lote completado en $\leq$10 minutos, una vez al d\'ia.
\item RNF-IA-08: cobertura $\geq$90\% de las aulas con sensor de ocupaci\'on activo.
\end{itemize}

\textbf{(d) Requisitos \'eticos y de privacidad.}
\begin{itemize}[nosep]
\item RNF-IA-09 (equidad): diferencia de detecci\'on $\leq$15 puntos porcentuales entre aulas con y sin sensor de consumo real.
\item RNF-IA-10 (equidad): los reportes no atribuyen responsabilidad individual a un docente por consumo que no controla directamente (climatizaci\'on autom\'atica de RF-13); se reportan agregados por aula y franja.
\item RNF-IA-11 (explicabilidad): cada anomal\'ia indica la variable que m\'as contribuy\'o, sin requerir consulta t\'ecnica adicional.
\item Finalidad: eficiencia energ\'etica institucional, no evaluaci\'on individual. Sin identidad de usuario/docente. Ninguna acci\'on correctiva autom\'atica.
\end{itemize}

\textbf{(e) Plan de monitoreo.} Indicadores: falsos positivos mensuales, cobertura de aulas analizadas. Umbral de deriva: falsos positivos $>$25\% durante dos meses consecutivos, o cobertura $<$80\%. Reentrenamiento: recalcular l\'inea base al inicio de cada periodo acad\'emico.


\subsection{RNF-17: Seguridad f\'isica (Safety)}
\textbf{Nota de numeraci\'on:} el ERS ya usa el identificador RNF-16 (Disponibilidad en modo degradado sin conectividad, Fiabilidad, sustentado en EV-15). Este requisito nuevo se numera \textbf{RNF-17} para no colisionar.

\textbf{Descripci\'on cuantificada:} antes de ejecutar un apagado autom\'atico de climatizaci\'on (RF-13, RF-16), el sistema debe verificar que el aula est\'a efectivamente desocupada mediante al menos dos lecturas de sensor de presencia consecutivas, separadas por $\geq$30~s, y debe registrar en bit\'acora (RF-23) el estado de ocupaci\'on que motiv\'o cada apagado. Ante una lectura inconsistente o ausente, el sistema \textbf{no} debe ejecutar el apagado y debe generar una alerta de ``verificaci\'on fallida'' (fallback seguro).

\textbf{Caracter\'istica ISO/IEC 25010:2023:} Seguridad f\'isica (Safety).

\textbf{Criterio de verificaci\'on:} prueba de inyecci\'on de lecturas inconsistentes/ausentes durante un ciclo de apagado autom\'atico $\rightarrow$ el sistema no ejecuta el comando y genera la alerta en el 100\% de un m\'inimo de 20 casos de prueba.

\textbf{Estado de validaci\'on:} derivaci\'on directa de RD-04 y RF-13/RF-16: \textbf{borrador PE5, pendiente de validaci\'on de campo}; no existe evidencia de entrevista porque el guion aplicado en la Entrega 3 (2A) no incluy\'o preguntas sobre fallas de sensor durante apagados autom\'aticos.

\textbf{Trazas:} RF-13, RF-16, RD-04, RD-13.


\subsection{Comparaci\'on de los dos componentes}

Los dos componentes comparten el marco de especificaci\'on pero difieren en casi todo lo dem\'as, y esas diferencias tienen consecuencias de dise\~no que conviene hacer expl\'icitas (Tabla~\ref{tab:comparacion-ia}).

\begin{table}[H]
\centering
\small
\begin{tabular}{|p{3.1cm}|p{5.1cm}|p{5.1cm}|}
\hline
\cellcolor{sigagreen}\textcolor{white}{\textbf{Dimensi\'on}} & \cellcolor{sigagreen}\textcolor{white}{\textbf{IA-01 Predictor de fallas}} & \cellcolor{sigagreen}\textcolor{white}{\textbf{IA-02 Consumo energ\'etico}} \\ \hline
Tipo de aprendizaje & Supervisado, clasificaci\'on binaria & No supervisado, detecci\'on de anomal\'ias \\ \hline
Etiquetado & Autom\'atico, desde tickets correctivos cerrados & No requiere etiquetado \\ \hline
Historial m\'inimo & 90 d\'ias por equipo & 60 d\'ias por aula \\ \hline
Frecuencia & Diaria, con latencia $\leq$2 s por inferencia & Diaria, por lote, en $\leq$10 minutos \\ \hline
M\'etrica principal & F1 $\geq$0,70 en holdout temporal & Falsos positivos $\leq$15\,\% \\ \hline
Consumidor & Infraestructura y Mantenimiento & Autoridades y Personal Administrativo \\ \hline
Consecuencia de un error & Una alerta informativa de m\'as o de menos & Una l\'inea de m\'as o de menos en un reporte \\ \hline
Riesgo de equidad & Ignorar equipos de aulas poco usadas & Penalizar aulas con menos instrumentaci\'on \\ \hline
\end{tabular}
\caption{Comparaci\'on de los dos componentes de inteligencia artificial}
\label{tab:comparacion-ia}
\end{table}

La diferencia con mayor efecto sobre la especificaci\'on es la del etiquetado. IA-01 puede entrenarse sin trabajo manual porque el propio sistema genera su variable objetivo: un ticket correctivo cerrado con causa de falla \emph{es} la etiqueta. Esa comodidad tiene un costo que se declara en la ficha como sesgo conocido: el historial de fallas depende de cu\'an diligente fue el registro manual antes de que el sistema existiera, de modo que el modelo aprender\'a tambi\'en los h\'abitos de registro del personal, no solo el comportamiento de los equipos.

IA-02, en cambio, no necesita etiquetas pero arrastra un problema distinto: en las aulas sin sensor de consumo dedicado trabaja sobre un indicador indirecto (ocupaci\'on m\'as estado de equipos) que es menos preciso. De ah\'i que RF-IA-04 obligue a \textbf{declarar en el propio resultado} si la l\'inea base proviene de medici\'on real o de indicador indirecto: sin ese dato, un reporte mezclar\'ia dos calidades de evidencia sin que el lector pueda distinguirlas.

\subsection{Plan de monitoreo y criterio de retirada}

Ninguno de los dos componentes se declara terminado al desplegarse. Ambos tienen indicadores en producci\'on, frecuencia de medici\'on, umbral de deriva y criterio de reentrenamiento, que la Tabla~\ref{tab:monitoreo-ia} resume.

\begin{table}[H]
\centering
\small
\begin{tabular}{|p{2.9cm}|p{5.2cm}|p{5.2cm}|}
\hline
\cellcolor{sigagreen}\textcolor{white}{\textbf{Elemento}} & \cellcolor{sigagreen}\textcolor{white}{\textbf{IA-01}} & \cellcolor{sigagreen}\textcolor{white}{\textbf{IA-02}} \\ \hline
Indicadores & F1 mensual, falsos negativos por segmento, latencia p95 & Falsos positivos mensuales, cobertura de aulas analizadas \\ \hline
Umbral de deriva & Ca\'ida de F1 mayor a 15 puntos respecto a la l\'inea base & Falsos positivos por encima del 25\,\% durante dos meses, o cobertura bajo el 80\,\% \\ \hline
Reentrenamiento & Cada 90 d\'ias o al activarse el umbral & Al inicio de cada periodo acad\'emico \\ \hline
Retirada & Autom\'atica si F1 baja de 0,60 en dos evaluaciones consecutivas & No aplica: el resultado solo alimenta un reporte revisado por humanos \\ \hline
\end{tabular}
\caption{Plan de monitoreo de los componentes de inteligencia artificial}
\label{tab:monitoreo-ia}
\end{table}

La asimetr\'ia de la \'ultima fila es deliberada. IA-01 se retira solo porque una predicci\'on degradada genera alertas que consumen tiempo de personal t\'ecnico; dejarlo operando con F1 bajo tiene un costo real. IA-02 no necesita retirada autom\'atica porque su salida nunca llega a un actuador ni a una persona sin pasar antes por la revisi\'on de quien firma el reporte.

\subsection{Nota sobre el uso de modelos de lenguaje en el propio proceso}

Conviene distinguir dos usos de la inteligencia artificial en este proyecto, porque se eval\'uan de forma distinta. Uno es el de los componentes IA-01 e IA-02, que forman parte del sistema especificado y son objeto de esta secci\'on. El otro es el uso de modelos de lenguaje como asistente del propio proceso de ingenier\'ia de requisitos, que es lo que estudia el componente emp\'irico del proyecto y lo que se declara en el Anexo E. Sobre ese segundo uso existe hoy evidencia sistematizada acerca de su alcance y sus l\'imites \cite{cheng2026genai,arora2024advancing}, gu\'ias metodol\'ogicas para aplicarlo a tareas de procesamiento de lenguaje natural en requisitos \cite{vogelsang2024using}, y estudios espec\'ificos sobre elicitaci\'on asistida \cite{ronanki2023investigating,gorer2023generating}, aseguramiento de la calidad de requisitos \cite{lubos2024leveraging} y estandarizaci\'on de su redacci\'on \cite{ray2023agile}. Ninguno de esos trabajos sostiene que el modelo sustituya el an\'alisis humano, y este informe no lo asume.

\subsection{S\'intesis}

Ambos componentes se clasifican como riesgo m\'inimo/limitado bajo el enfoque del Reglamento (UE) 2024/1689, usado como marco de referencia de buenas pr\'acticas: no como obligaci\'on legal directa, dado que el sistema opera \'unicamente en Ecuador sin nexo con la Uni\'on Europea. La obligaci\'on legal real y aplicable es la Ley Org\'anica de Protecci\'on de Datos Personales del Ecuador \cite{asamblea2021lopdp}, cuyos art\'iculos 13 y 14 originan RF-24 y RF-25 y se mapean requisito por requisito en la \S3.4 del ERS. Ninguno de los dos componentes ejecuta una acci\'on autom\'atica irreversible sobre personas o equipos sin supervisi\'on humana: toda predicci\'on de alto riesgo o anomal\'ia de consumo pasa por revisi\'on de Infraestructura/Autoridades antes de generar un ticket o una decisi\'on de inversi\'on.
